% Generated from tex/chapters/ch04/05_ソフトウェア構築の基礎.tex for the SWEBOK structure tree
\subsubsection{資産の再利用}

再利用とは、既存の資産を別の問題解決に使うことである。構築で再利用される資産には、フレームワーク、ライブラリ、モジュール、コンポーネント、ソースコード、商用既製品、クラウドサービスなどがある。

再利用には二つの側面がある。第一は、再利用されることを意識して作る「再利用のための構築」である。第二は、既存の資産を使って新しい解決策を作る「再利用による構築」である。前者では、汎用性、可変性、文書化、テストが重要になる。後者では、選定、評価、統合、ライセンス、保守、脆弱性管理が重要になる。

\begin{keybox}{注意}
再利用は必ず得とは限らない。再利用した部品に脆弱性、ライセンス問題、保守停止、過剰な依存があれば、後から大きな費用が発生する。再利用するものは、品質と責任範囲を確認する必要がある。
\end{keybox}
